% !TeX program = lualatex
% !TeX root = o predmetu.tex
% ============================================================
% Zajednički preamble za sva predavanja RWA kolegija
% Uključiti s: % ============================================================
% Zajednički preamble za sva predavanja RWA kolegija
% Uključiti s: % ============================================================
% Zajednički preamble za sva predavanja RWA kolegija
% Uključiti s: \input{../common/preamble.tex}
% ============================================================

\documentclass[aspectratio=149,11pt]{beamer}
\usetheme{CambridgeUS}

\usepackage{fontspec}              % radi samo sa XeLaTeX/LuaLaTeX
\setsansfont{Latin Modern Sans}    % sans-serif font
\setmonofont{Latin Modern Mono}    % monospace font za kod
\usepackage[croatian]{babel}
\usepackage{csquotes}              % \enquote{} → hrvatski navodnici „..."
\usepackage{lmodern}
\usepackage{booktabs}
\usepackage{listings}
\usepackage{xcolor}
\usepackage{hyperref}

% ---------- Autor i logo ----------
\author{mag.ing. Josip Torić (SIT UNIZD)}
\titlegraphic{\includegraphics[height=2.5cm]{../common/SIT_logo_tekst1.pdf}}

% ---------- Listings ----------
\definecolor{codebg}{RGB}{248,248,248}
\definecolor{codeframeclr}{RGB}{220,220,220}
\lstdefinelanguage{bash}{
  morekeywords={sudo,apt,brew,pip,pipx,uv,winget,python,git,flask,source,deactivate,code,gunicorn,docker,compose,heroku,render,railway,node,npm,pnpm,yarn,uvicorn,alembic,pytest,ruff,black,precommit},
  sensitive=false,
  morecomment=[l]{\#}
}
\lstdefinelanguage{python}{
  morekeywords={from,import,def,class,return,if,elif,else,for,in,while,try,except,raise,with,as,not,and,or,is,lambda,yield,True,False,None,await,async},
  sensitive=true,
  morecomment=[l]{\#},
  morestring=[b]',
  morestring=[b]"
}
\lstdefinelanguage{sql}{
  morekeywords={SELECT,FROM,WHERE,INSERT,INTO,VALUES,UPDATE,SET,DELETE,CREATE,TABLE,PRIMARY,KEY,FOREIGN,REFERENCES,UNIQUE,NOT,NULL,INDEX,ON,ALTER,DROP,CONSTRAINT},
  sensitive=false,
  morecomment=[l]{--}
}
\lstdefinelanguage{json}{
  morecomment=[l]{//},
  morestring=[b]",
  showstringspaces=false
}
\lstdefinelanguage{typescript}{
  morekeywords={import,export,from,const,let,var,function,return,if,else,for,of,in,while,switch,case,break,async,await,try,catch,finally,throw,new,class,this,extends,implements,interface,type,enum,readonly,public,private,protected,static,get,set,true,false,null,undefined,void,string,number,boolean,any,unknown,never,typeof,keyof,as},
  sensitive=true,
  morecomment=[l]{//},
  morecomment=[s]{/*}{*/},
  morestring=[b]',
  morestring=[b]",
  morestring=[b]`
}
\lstalias{javascript}{typescript}
\lstset{
  backgroundcolor=\color{codebg},
  basicstyle=\ttfamily\footnotesize,
  frame=single,
  rulecolor=\color{codeframeclr},
  frameround=tttt,
  breaklines=true,
  tabsize=2,
  showstringspaces=false,
  upquote=true,
  columns=fullflexible
}

% ---------- Helper macros ----------
\newcommand{\codefile}[1]{\texttt{#1}}
\newcommand{\kbd}[1]{\texttt{#1}}

% ---------- Code frame environment ----------
% Automatski dodaje fragile i shrink=5, koristi se umjesto \begin{frame}
% Primjer: \begin{codeframe}{Naslov} ... \end{codeframe}
% Primjer s opcijama: \begin{codeframe}[shrink=15]{Naslov} ... \end{codeframe}
\newenvironment{codeframe}[2][]
  {\begin{frame}[fragile,shrink=5,#1]{#2}}
  {\end{frame}}

% ============================================================

\documentclass[aspectratio=149,11pt]{beamer}
\usetheme{CambridgeUS}

\usepackage{fontspec}              % radi samo sa XeLaTeX/LuaLaTeX
\setsansfont{Latin Modern Sans}    % sans-serif font
\setmonofont{Latin Modern Mono}    % monospace font za kod
\usepackage[croatian]{babel}
\usepackage{csquotes}              % \enquote{} → hrvatski navodnici „..."
\usepackage{lmodern}
\usepackage{booktabs}
\usepackage{listings}
\usepackage{xcolor}
\usepackage{hyperref}

% ---------- Autor i logo ----------
\author{mag.ing. Josip Torić (SIT UNIZD)}
\titlegraphic{\includegraphics[height=2.5cm]{../common/SIT_logo_tekst1.pdf}}

% ---------- Listings ----------
\definecolor{codebg}{RGB}{248,248,248}
\definecolor{codeframeclr}{RGB}{220,220,220}
\lstdefinelanguage{bash}{
  morekeywords={sudo,apt,brew,pip,pipx,uv,winget,python,git,flask,source,deactivate,code,gunicorn,docker,compose,heroku,render,railway,node,npm,pnpm,yarn,uvicorn,alembic,pytest,ruff,black,precommit},
  sensitive=false,
  morecomment=[l]{\#}
}
\lstdefinelanguage{python}{
  morekeywords={from,import,def,class,return,if,elif,else,for,in,while,try,except,raise,with,as,not,and,or,is,lambda,yield,True,False,None,await,async},
  sensitive=true,
  morecomment=[l]{\#},
  morestring=[b]',
  morestring=[b]"
}
\lstdefinelanguage{sql}{
  morekeywords={SELECT,FROM,WHERE,INSERT,INTO,VALUES,UPDATE,SET,DELETE,CREATE,TABLE,PRIMARY,KEY,FOREIGN,REFERENCES,UNIQUE,NOT,NULL,INDEX,ON,ALTER,DROP,CONSTRAINT},
  sensitive=false,
  morecomment=[l]{--}
}
\lstdefinelanguage{json}{
  morecomment=[l]{//},
  morestring=[b]",
  showstringspaces=false
}
\lstdefinelanguage{typescript}{
  morekeywords={import,export,from,const,let,var,function,return,if,else,for,of,in,while,switch,case,break,async,await,try,catch,finally,throw,new,class,this,extends,implements,interface,type,enum,readonly,public,private,protected,static,get,set,true,false,null,undefined,void,string,number,boolean,any,unknown,never,typeof,keyof,as},
  sensitive=true,
  morecomment=[l]{//},
  morecomment=[s]{/*}{*/},
  morestring=[b]',
  morestring=[b]",
  morestring=[b]`
}
\lstalias{javascript}{typescript}
\lstset{
  backgroundcolor=\color{codebg},
  basicstyle=\ttfamily\footnotesize,
  frame=single,
  rulecolor=\color{codeframeclr},
  frameround=tttt,
  breaklines=true,
  tabsize=2,
  showstringspaces=false,
  upquote=true,
  columns=fullflexible
}

% ---------- Helper macros ----------
\newcommand{\codefile}[1]{\texttt{#1}}
\newcommand{\kbd}[1]{\texttt{#1}}

% ---------- Code frame environment ----------
% Automatski dodaje fragile i shrink=5, koristi se umjesto \begin{frame}
% Primjer: \begin{codeframe}{Naslov} ... \end{codeframe}
% Primjer s opcijama: \begin{codeframe}[shrink=15]{Naslov} ... \end{codeframe}
\newenvironment{codeframe}[2][]
  {\begin{frame}[fragile,shrink=5,#1]{#2}}
  {\end{frame}}

% ============================================================

\documentclass[aspectratio=149,11pt]{beamer}
\usetheme{CambridgeUS}

\usepackage{fontspec}              % radi samo sa XeLaTeX/LuaLaTeX
\setsansfont{Latin Modern Sans}    % sans-serif font
\setmonofont{Latin Modern Mono}    % monospace font za kod
\usepackage[croatian]{babel}
\usepackage{csquotes}              % \enquote{} → hrvatski navodnici „..."
\usepackage{lmodern}
\usepackage{booktabs}
\usepackage{listings}
\usepackage{xcolor}
\usepackage{hyperref}

% ---------- Autor i logo ----------
\author{mag.ing. Josip Torić (SIT UNIZD)}
\titlegraphic{\includegraphics[height=2.5cm]{../common/SIT_logo_tekst1.pdf}}

% ---------- Listings ----------
\definecolor{codebg}{RGB}{248,248,248}
\definecolor{codeframeclr}{RGB}{220,220,220}
\lstdefinelanguage{bash}{
  morekeywords={sudo,apt,brew,pip,pipx,uv,winget,python,git,flask,source,deactivate,code,gunicorn,docker,compose,heroku,render,railway,node,npm,pnpm,yarn,uvicorn,alembic,pytest,ruff,black,precommit},
  sensitive=false,
  morecomment=[l]{\#}
}
\lstdefinelanguage{python}{
  morekeywords={from,import,def,class,return,if,elif,else,for,in,while,try,except,raise,with,as,not,and,or,is,lambda,yield,True,False,None,await,async},
  sensitive=true,
  morecomment=[l]{\#},
  morestring=[b]',
  morestring=[b]"
}
\lstdefinelanguage{sql}{
  morekeywords={SELECT,FROM,WHERE,INSERT,INTO,VALUES,UPDATE,SET,DELETE,CREATE,TABLE,PRIMARY,KEY,FOREIGN,REFERENCES,UNIQUE,NOT,NULL,INDEX,ON,ALTER,DROP,CONSTRAINT},
  sensitive=false,
  morecomment=[l]{--}
}
\lstdefinelanguage{json}{
  morecomment=[l]{//},
  morestring=[b]",
  showstringspaces=false
}
\lstdefinelanguage{typescript}{
  morekeywords={import,export,from,const,let,var,function,return,if,else,for,of,in,while,switch,case,break,async,await,try,catch,finally,throw,new,class,this,extends,implements,interface,type,enum,readonly,public,private,protected,static,get,set,true,false,null,undefined,void,string,number,boolean,any,unknown,never,typeof,keyof,as},
  sensitive=true,
  morecomment=[l]{//},
  morecomment=[s]{/*}{*/},
  morestring=[b]',
  morestring=[b]",
  morestring=[b]`
}
\lstalias{javascript}{typescript}
\lstset{
  backgroundcolor=\color{codebg},
  basicstyle=\ttfamily\footnotesize,
  frame=single,
  rulecolor=\color{codeframeclr},
  frameround=tttt,
  breaklines=true,
  tabsize=2,
  showstringspaces=false,
  upquote=true,
  columns=fullflexible
}

% ---------- Helper macros ----------
\newcommand{\codefile}[1]{\texttt{#1}}
\newcommand{\kbd}[1]{\texttt{#1}}

% ---------- Code frame environment ----------
% Automatski dodaje fragile i shrink=5, koristi se umjesto \begin{frame}
% Primjer: \begin{codeframe}{Naslov} ... \end{codeframe}
% Primjer s opcijama: \begin{codeframe}[shrink=15]{Naslov} ... \end{codeframe}
\newenvironment{codeframe}[2][]
  {\begin{frame}[fragile,shrink=5,#1]{#2}}
  {\end{frame}}


\title[Razvoj web aplikacija]{Razvoj web aplikacija - Predavanje 0 \\ O predmetu}
\date{05.10.2026.}

\begin{document}

\frame{\titlepage}

% ------------------------------
\begin{frame}{Plan predavanja}
\begin{enumerate}
  \item Što je kolegij i koji je cilj
  \item Što gradimo: sustav prijava na powerlifting natjecanje
  \item Organizacija nastave, pravila, bodovanje i obrane
  \item Tehnologije, arhitektura i razvojni workflow
  \item Kako učiti ovaj kolegij
\end{enumerate}

\vspace{0.6em}
Ne radimo \enquote{tutorial app}. Radimo \textbf{inženjerski} projekt: pravila, sigurnost, testovi, deploy.
\end{frame}

% ------------------------------
\begin{frame}{Cilj kolegija}
\begin{itemize}
  \item Izraditi \textbf{full-stack} web aplikaciju: backend + frontend + baza
  \item Razumjeti \textbf{arhitekturu} web sustava i trade-offe
  \item Ovladati \textbf{API dizajnom} i modeliranjem podataka
  \item Implementirati \textbf{autentikaciju} (JWT) i \textbf{autorizaciju} (role/ownership)
  \item Izgraditi SPA frontend: \textbf{Vue 3}, routing, state management, UX discipline
  \item Uspostaviti \textbf{workflow}: Git, commitovi, testiranje, CI i deploy
\end{itemize}

\vspace{0.4em}
\textbf{Ishod:} student može samostalno izraditi, testirati i deployati manju web aplikaciju.
\end{frame}

% ------------------------------
\begin{frame}{Što očekujem od vas}
\begin{itemize}
  \item Redovito praćenje predavanja i razumijevanje koncepata, ne samo kopiranje
  \item Rad u repozitoriju: \textbf{commit discipline} i čitljiv kod
  \item Postavljanje pitanja i rješavanje problema - debugging kao vještina
  \item Minimalna higijena: formatiranje, error handling, logika na pravom mjestu
\end{itemize}

\vspace{0.4em}
Najveća razlika između prolaza i pada je: \textbf{razumijete li zašto kod radi}.
\end{frame}

% ------------------------------
\begin{frame}{Projekt: Sustav prijava na powerlifting natjecanje}
\begin{columns}[T]
\column{0.48\textwidth}
\textbf{Uloge}
\begin{itemize}
  \item \textbf{Admin}: kreira natjecanje, kreira klubove i login podatke, vidi sve prijave
  \item \textbf{Klub}: registrira natjecatelje i prijavljuje ih na natjecanje
\end{itemize}

\column{0.48\textwidth}
\textbf{Temeljna domena}
\begin{itemize}
  \item Klub, Korisnik, Natjecatelj, Natjecanje, Prijava
  \item Rokovi nominacija: \textbf{prelim} i \textbf{final}
\end{itemize}
\end{columns}
\end{frame}

% ------------------------------
\begin{frame}{Pravila rokova - nominacije}
\begin{tabular}{@{}lp{8cm}@{}}
  \toprule
  \textbf{Faza} & \textbf{Što klub može} \\
  \midrule
  Do \textbf{prelim} roka &
    Dodavati nove prijave, uređivati podatke prijave, odjaviti prijavu \\
  \midrule
  Između \textbf{prelim} i \textbf{final} &
    Odjaviti natjecatelja ili promijeniti kategoriju - \textbf{ne može} dodavati nove prijave \\
  \midrule
  Nakon \textbf{final} roka &
    Nema izmjena \\
  \bottomrule
\end{tabular}

\vspace{0.6em}
Natjecanje ima dva ključna roka: \textbf{preliminarne nominacije} (\codefile{prelim\_deadline}) i \textbf{finalne nominacije} (\codefile{final\_deadline}).
\end{frame}

% ------------------------------
\begin{frame}{Zašto je ova tema dobra za učenje?}
\begin{itemize}
  \item Prirodno sadrži \textbf{role} i \textbf{ownership} - admin vidi sve, klub samo svoje
  \item Ima \textbf{workflow} i vremenska poslovna pravila - realan backend, ne samo CRUD
  \item Relacije i constrainti - realan rad s bazom
  \item Frontend je zanimljiviji: tablice, filteri, statusi, onemogućeni gumbi po fazi
  \item Obrana po kodu je jasna: može se postaviti \textbf{precizno pitanje} i provjeriti razumijevanje
\end{itemize}
\end{frame}

% ------------------------------
\begin{frame}{Tehnologije koje koristimo}
\begin{block}{Backend}
FastAPI + Uvicorn \;|\; Pydantic \;|\; SQLAlchemy 2.0 \;|\; Alembic \;|\; PostgreSQL
\end{block}

\begin{block}{Frontend}
Vue 3 + Vite + TypeScript \;|\; Vue Router \;|\; Pinia \;|\; Axios
\end{block}

\begin{block}{Ops / workflow}
Docker Compose \;|\; Git \;|\; uv \;|\; pytest \;|\; ruff \;|\; CI \;|\; Deploy (Railway ili Render)
\end{block}
\end{frame}

% ------------------------------
\begin{frame}{Arhitektura aplikacije}

\textbf{Frontend (SPA)} komunicira s backendom preko \textbf{REST API-ja}. Backend implementira:
\begin{itemize}
  \item autentikaciju (JWT) i autorizaciju (role/ownership)
  \item poslovna pravila i rad s bazom
  \item dokumentaciju API-ja (OpenAPI/Swagger)
\end{itemize}

\vspace{0.4em}
Baza je \textbf{jedini izvor istine} za podatke.

\vspace{0.4em}
Frontend je \enquote{glup} za pravila - backend je autoritet. UI može sakriti opcije, ali backend \textbf{mora zabraniti}.
\end{frame}

% ------------------------------
\begin{frame}{Struktura projekta}
\begin{columns}[T]
\column{0.48\textwidth}
\textbf{Backend}
\begin{itemize}
  \item \codefile{routers/} - HTTP rute i parametri
  \item \codefile{schemas/} - Pydantic DTO modeli
  \item \codefile{models/} - SQLAlchemy modeli
  \item \codefile{services/} - poslovna logika
  \item \codefile{repositories/} - DB upiti
  \item \codefile{core/} - config, auth, errors
\end{itemize}

\column{0.48\textwidth}
\textbf{Frontend}
\begin{itemize}
  \item \codefile{src/views/} - stranice
  \item \codefile{src/components/} - UI komponente
  \item \codefile{src/stores/} - Pinia storeovi
  \item \codefile{src/router/} - rute i guardovi
  \item \codefile{src/api/} - API client wrapper
\end{itemize}
\end{columns}

\vspace{0.4em}
Cilj je izbjeći \enquote{sve u jednom file-u}. Svaki sloj ima jednu odgovornost.
\end{frame}

% ------------------------------
\begin{frame}[shrink=15]{Raspored predavanja po temama}
	\begin{tabular}{@{}lll@{}}
		\toprule
		\textbf{Pred.} & \textbf{Datum} & \textbf{Tema} \\
		\midrule
		0  & 05.10. & O predmetu \\
		1  & 05.10. & Backend arhitektura i skeleton projekta \\
		2  & 12.10. & Modeliranje podataka i rad s bazom \\
		3  & 19.10. & Autentikacija i sigurnost (JWT) \\
		4  & 26.10. & Autorizacija, role i ownership \\
		5  & 02.11. & API dizajn, validacija i CRUD domene \\
		6  & 09.11. & Poslovna pravila i workflow u backendu \\
		-  & 16.11. & \textbf{Obrana 1. dijela projekta (backend)} \\
		7  & 23.11. & Arhitektura frontend aplikacija (Vue) \\
		8  & 30.11. & Komunikacija frontend-backend i auth flow \\
		9  & 07.12. & State management i UX discipline \\
		10 & 14.12. & Kompleksni UI, tablice i administracija \\
		11 & 21.12. & Testiranje full-stack aplikacije \\
		& 28.12. - 04.01. & \emph{Božićni praznici} \\
		12 & 11.01. & Frontend u produkciji i deploy \\
		-  & 18.01. & \textbf{Obrana 2. dijela projekta (full-stack)} \\
		\bottomrule
	\end{tabular}
\end{frame}

% ------------------------------
\begin{frame}{Projekt studenata - zahtjevi}
	\begin{itemize}
		\item Tijekom semestra izrađujete \textbf{vlastitu full-stack web aplikaciju}
		\item Aplikacija mora sadržavati:
		\begin{itemize}
			\item autentikaciju (JWT) i barem dvije role
			\item najmanje dva entiteta i njihove relacije
			\item ownership - korisnik vidi i uređuje svoje podatke
			\item poslovna pravila: workflow, statusi ili rokovi
			\item frontend SPA s osnovnim UX stanjima: loading, error, empty
		\end{itemize}
		\item \textbf{Predaja 1. dijela:} backend funkcionalan, predan na GitHub - 45 bodova
		\item \textbf{Predaja 2. dijela:} full-stack na GitHubu i \textbf{deployan} - 45 bodova
		\item Ako ne napravite sve tijekom nastave, imate još jedan termin predaje za vrijeme ispitnih rokova u veljači
	\end{itemize}
\end{frame}

% ------------------------------
\begin{frame}{Bodovanje}
Dolazak na predavanje nosi \textbf{1 bod}, najviše \textbf{10 bodova} u semestru.

\vspace{0.6em}
\begin{tabular}{@{}lr@{}}
  \toprule
  \textbf{Komponenta} & \textbf{Bodovi} \\
  \midrule
  1. dio projekta (backend) & 45 \\
  2. dio projekta (frontend + deploy) & 45 \\
  Dolasci na predavanja & 10 \\
  \midrule
  \textbf{Ukupno} & \textbf{100} \\
  \bottomrule
\end{tabular}

\vspace{0.6em}
Predavanja ima 12, a bodovanih dolazaka 10 - možete propustiti dva termina bez posljedica. \\
Dolazak nije obavezan, ali sada nosi dio ocjene. Najviše ćete ionako dobiti ako dolazite jer želite nešto naučiti.
\end{frame}

% ------------------------------
\begin{frame}{Primjeri tema projekata}
	\begin{itemize}
		\item Sustav rezervacija - termini, dvorane, treneri, konzultacije
		\item Event management - događaji, prijave, kapacitet, check-in
		\item Task / workflow sustav - projekti, zadaci, statusi, komentari
		\item Coach \& athlete tracker - programi, treninzi, napredak
		\item Dokument approval sustav - zahtjevi, odobravanje, verzije
		\item Mini marketplace - oglasi, poruke, favoriti
		\item Evidencija članstva - klubovi, članovi, uplate, statusi
	\end{itemize}

	\vspace{0.6em}
	Odaberite temu koja ima \textbf{uloge i pravila}, a ne samo jednostavan CRUD.
\end{frame}

% ------------------------------
\begin{frame}{Kako izgleda obrana}
Demonstracija funkcionalnosti, pregled repozitorija i \textbf{ključnih dijelova koda}, pitanja o odlukama i edge caseovima. Moguća mini promjena uživo: dodaj polje, migraciju, prilagodi DTO/UI.

\vspace{0.6em}
\begin{block}{Tipična pitanja}
\begin{itemize}
  \item Gdje provjeravaš ownership?
  \item Kako backend zabranjuje izmjene nakon roka?
  \item Zašto je status code 401/403/409/422?
  \item Gdje i kako hendlaš 401 na frontendu?
\end{itemize}
\end{block}
\end{frame}

% ------------------------------
\begin{frame}{Lokalni setup}
\begin{itemize}
  \item Git
  \item Python 3.11+ za backend
  \item Docker Desktop + Docker Compose za bazu
  \item Node.js 22+ za frontend
  \item Editor: VS Code / IntelliJ / drugi
\end{itemize}
\end{frame}

% ------------------------------
\begin{frame}[fragile,shrink=5]{Pokretanje projekta lokalno}
\begin{lstlisting}[language=bash]
# 1) Baza (Postgres)
docker compose up -d db

# 2) Backend
cd api
uv venv
source .venv/bin/activate
uv pip install -r requirements.txt
alembic upgrade head
python -m app.seed
uvicorn app.main:app --reload

# 3) Frontend
cd ../web
npm install
npm run dev
\end{lstlisting}
\end{frame}

% ------------------------------
\begin{frame}{Kvaliteta koda: što je \enquote{dovoljno dobro}}
\begin{columns}[T]
\column{0.48\textwidth}
\textbf{Backend}
\begin{itemize}
  \item Validacija inputa (Pydantic)
  \item Dosljedni errori (standardiziran JSON)
  \item Logika u service layeru, ne u routeru
  \item Testovi za ključna pravila
\end{itemize}

\column{0.48\textwidth}
\textbf{Frontend}
\begin{itemize}
  \item Loading/error/empty stanja
  \item Route guardovi za role
  \item Form validation
  \item Bez hardkodiranih URL-ova (env)
\end{itemize}
\end{columns}
\end{frame}

% ------------------------------
\begin{frame}{Testiranje}
Test nije \enquote{još jedan zadatak} - to je \textbf{osiguranje} da pravila vrijede. Fokus testova u ovom projektu:

\begin{itemize}
  \item Autentikacija - login/refresh
  \item Autorizacija - role/ownership
  \item Rokovi - prije prelim / između / nakon final
  \item Unikatan constraint - nema duple prijave
\end{itemize}

\vspace{0.4em}
Na obrani ću vas pitati \textbf{što točno test pokriva} i \textbf{zašto je to bitno}.
\end{frame}

% ------------------------------
\begin{frame}{Deploy: što znači \enquote{radi u produkciji}}
Deploy nije samo \enquote{imam URL}. Morate razumjeti:

\begin{itemize}
  \item Env varovi i tajne - JWT secret, DB URL
  \item CORS - frontend domena $\leftrightarrow$ backend
  \item Migracije u produkciji - alembic upgrade
  \item Error logging - što logiramo, što ne
\end{itemize}

\vspace{0.4em}
Koristit ćemo moderni PaaS - Railway ili Render.
\end{frame}

% ------------------------------
\begin{frame}{Akademska čestitost i AI alati}
AI možete koristiti za pomoć, ali na obrani morate \textbf{razumjeti i objasniti svoj kod}. Kopirani kod bez razumijevanja se vidi odmah na usmenom.

\vspace{0.4em}
Ako koristite vanjske snippete, morate znati:
\begin{itemize}
  \item Gdje su rizici za sigurnost
  \item Što se događa u edge caseovima
  \item Kako to testirate
\end{itemize}
\end{frame}

% ------------------------------
\begin{frame}{Kako učiti ovaj kolegij}
\begin{itemize}
  \item Svako predavanje: pogledaj diffove, reproduciraj korake u svom repou
  \item Ne gomilaj dugačke sesije - radi \textbf{malo i često}
  \item Debugging rutina:
    \begin{itemize}
      \item pročitaj error, izoliraj uzrok, napravi najmanji repro
      \item provjeri request/response - status code, payload
      \item koristi logove i devtools
    \end{itemize}
  \item Kad zapne: pokaži \textbf{konkretan problem} - log, request, commit
\end{itemize}
\end{frame}
\end{document}
